Como el problema general de determinar si dos programas son equivalentes es indecidible, no existe un procedimiento universal capaz de resolver esta cuestión para cualquier par de programas. En consecuencia, el estudio de la equivalencia debe abordarse mediante criterios de indistinguibilidad establecidos por aspectos particulares del comportamiento computacional. 

Entre los enfoques más relevantes se encuentran la \emph{equivalencia operacional}, la \emph{equivalencia denotativa} y la \emph{equivalencia axiomática}, cada una asociada a una forma distinta de describir y razonar sobre el significado de los programas.

A través de la formalización de estas tres nociones de equivalencia para el lenguaje \textsc{Comm} mediante el asistente de pruebas \textsc{Rocq}, este trabajo permitió constatar que una misma propiedad puede requerir estrategias de demostración significativamente distintas según el marco semántico adoptado. Mientras que las demostraciones operacionales exigen reconstruir derivaciones y analizar explícitamente las transformaciones de estado producidas durante la ejecución, las demostraciones denotativas se apoyan en propiedades matemáticas de las funciones que representan el significado de los programas, como la extensionalidad funcional. Por su parte, la equivalencia axiomática basada en teoría de modelos mostró cómo el razonamiento sobre propiedades lógicas de los programas puede fundamentarse en la estructura del modelo de estados subyacente.

En conjunto, los resultados obtenidos muestran que las distintas nociones de equivalencia no deben entenderse como conceptos rivales, sino como perspectivas complementarias para estudiar el comportamiento de los programas. De manera general, dos programas pueden considerarse equivalentes cuando resultan indistinguibles respecto de los criterios observables establecidos por una determinada semántica: en la semántica operacional, la equivalencia se establece a partir del comportamiento de ejecución y de las transformaciones de estado producidas por los programas; en la semántica denotativa, se determina mediante la igualdad de las funciones que representan su significado matemático; y en la semántica axiomática, se fundamenta en la imposibilidad de distinguir los programas mediante las propiedades lógicas que satisfacen.

Comprender sus diferencias, alcances y relaciones proporciona una base teórica sólida para áreas como la verificación formal, el diseño de lenguajes de programación y la optimización segura de compiladores. En estos contextos, la capacidad de demostrar rigurosamente que una transformación preserva el significado de un programa constituye un requisito fundamental para garantizar la corrección y confiabilidad de los sistemas de software.

\subsection{Trabajo a futuro}

Si bien este trabajo permitió formalizar y comparar distintas nociones de equivalencia de programas para el lenguaje \textsc{Comm}, permanecen abiertos diversos problemas cuya exploración podría ampliar tanto el alcance teórico como el grado de formalización de los resultados obtenidos.

En primer lugar, el estudio se centró exclusivamente en propiedades relacionadas con la semántica dinámica de los programas. No obstante, en la Sección~\ref{sec:equivprogram} se introdujo la noción de equivalencia de tipos, la cual pertenece al ámbito de la semántica estática. Una posible extensión consiste en analizar las distintas formas de equivalencia asociadas a los sistemas de tipos y estudiar las relaciones existentes entre las nociones de equivalencia estática y dinámica dentro de un mismo lenguaje de programación.

Por otra parte, la formalización de la semántica denotativa de $\textsc{Comm}$ se restringió a programas sin constructores iterativos. En consecuencia, resulta de interés extender la implementación en \textsc{Rocq} para incorporar la instrucción \textit{while}. Esta extensión requiere el estudio de herramientas matemáticas más avanzadas, particularmente la teoría de dominios y la teoría de puntos fijos, las cuales permiten modelar formalmente el significado de las construcciones recursivas e iterativas.

Finalmente, el análisis de la semántica axiomática reveló la existencia de dos nociones de equivalencia: una basada en la teoría de modelos y otra basada en la teoría de la prueba. En este trabajo únicamente se desarrolló la primera de ellas. Una línea natural de investigación consiste en formalizar la equivalencia axiomática desde la perspectiva de la teoría de la prueba y estudiar su relación con la equivalencia definida mediante teoría de modelos. En particular, sería necesario analizar las propiedades de corrección y completitud del sistema lógico utilizado para expresar propiedades de los programas. Si dichas propiedades se verifican para el lenguaje $\textsc{Comm}$, podría plantearse una caracterización unificada de la equivalencia axiomática que permita demostrar la coincidencia entre ambas perspectivas.